\documentclass[11pt,a4paper]{article}

% ----------------------------------------------------------------------------
% Pedagogical write-up of the DQMC algorithm implemented in qmc_dqmc.
% Compile with:  pdflatex dqmc.tex   (twice for the table of contents)
% ----------------------------------------------------------------------------

\usepackage[margin=1in]{geometry}
\usepackage[utf8]{inputenc}
\usepackage[T1]{fontenc}
\usepackage{lmodern}
\usepackage{microtype}
\usepackage{amsmath,amssymb,amsthm,mathtools}
\usepackage{physics}
\usepackage{braket}
\usepackage{bm}
\usepackage{tikz}
\usetikzlibrary{positioning,arrows.meta,decorations.pathreplacing,calc}
\usepackage{listings}
\usepackage{xcolor}
\usepackage{hyperref}
\usepackage{cleveref}
\usepackage{booktabs}
\usepackage{enumitem}

\hypersetup{
    colorlinks=true,
    linkcolor=blue!50!black,
    citecolor=blue!50!black,
    urlcolor=blue!50!black,
    pdftitle={Determinant Quantum Monte Carlo for the Hubbard model},
    pdfauthor={qmc\_sse contributors}
}

\renewcommand{\Tr}{\operatorname{Tr}}
\newcommand{\eg}{e.g.\ }
\newcommand{\ie}{i.e.\ }
\newcommand{\Hop}{\hat H}
\newcommand{\Bm}[2]{B_{#1,#2}}
\newcommand{\Ns}{N_{s}}
\newcommand{\Lt}{L_\tau}
\newcommand{\dt}{\Delta\tau}

\theoremstyle{plain}
\newtheorem{proposition}{Proposition}

\title{\bfseries Determinant Quantum Monte Carlo for the Hubbard model\\
       \large A self-contained derivation of the BSS algorithm and its
              implementation in \texttt{qmc\_dqmc}}
\author{\texttt{qmc\_dqmc} contributors}
\date{}

\begin{document}
\maketitle
\tableofcontents
\bigskip

\section{Overview}

The Determinant Quantum Monte Carlo (DQMC) algorithm of
Blankenbecler, Scalapino and Sugar~\cite{BSS1981}, equipped with the
Hirsch~\cite{Hirsch1983} discrete Hubbard--Stratonovich transformation,
is the canonical tool for unbiased finite-temperature simulations of
the repulsive Hubbard model. This document derives the algorithm
implemented in \texttt{qmc\_dqmc}, focusing on the practical equations
needed to read the code in \texttt{include/qmc/dqmc\_engine.hpp}.

The companion document \texttt{sse\_qmc.tex} treats the SSE algorithm
for spin systems; the present notes are independent and use a similar
notation. See also the implementation summary in \texttt{docs/DQMC.md}.

\section{Model and conventions}

We simulate the single-band Hubbard model on a generic lattice with
periodic boundary conditions:
\begin{equation}
\Hop = -t \sum_{\langle ij\rangle, \sigma}
            \left( c^{\dagger}_{i\sigma} c_{j\sigma} + \text{h.c.} \right)
       - \mu \sum_{i\sigma} n_{i\sigma}
       + U \sum_{i} \left(n_{i\uparrow} - \tfrac{1}{2}\right)
                    \left(n_{i\downarrow} - \tfrac{1}{2}\right) .
\label{eq:hubbard}
\end{equation}
The symmetric \(\bigl(n - \tfrac{1}{2}\bigr)\) form makes
\(\mu = 0\) coincide with half-filling on bipartite lattices by
particle--hole symmetry. We collect the kinetic and chemical-potential
terms into a single-particle matrix
\(K\in\mathbb{R}^{\Ns\times\Ns}\) with
\(K_{ij} = -t\) for nearest neighbours and \(K_{ii} = -\mu\).

\section{Trotter decomposition and HS transformation}

Discretise imaginary time as \(\beta = \Lt \dt\). To leading order
in \(\dt\),
\begin{equation}
e^{-\beta \Hop} = \prod_{l = 0}^{\Lt - 1}
   e^{-\dt K} \, e^{-\dt \, V} + \mathcal{O}(\dt^2),
\label{eq:trotter}
\end{equation}
where \(V = U\sum_i (n_{i\uparrow}-\tfrac12)(n_{i\downarrow}-\tfrac12)\).
Hirsch's discrete HS identity decouples the on-site quartic term:
\begin{equation}
e^{-\dt U (n_{i\uparrow}-\tfrac12)(n_{i\downarrow}-\tfrac12)}
   = \frac{e^{-\dt U/4}}{2}
     \sum_{s = \pm 1} e^{\alpha s (n_{i\uparrow} - n_{i\downarrow})},
\qquad \cosh \alpha = e^{\dt U / 2}.
\label{eq:HS}
\end{equation}
Introducing one Hubbard--Stratonovich field \(s_{l,i}\in\{-1,+1\}\) at
every site of every Trotter slice, we obtain
\begin{equation}
Z = \Tr e^{-\beta\Hop}
  = \frac{e^{-\beta U \Ns / 4}}{2^{\Ns \Lt}}
    \sum_{\{s_{l,i}\}}
    \det M_{\uparrow}(s) \, \det M_{\downarrow}(s),
\label{eq:Z}
\end{equation}
with the single-particle propagation matrices
\begin{equation}
M_\sigma(s) = \mathbb{1}
            + \Bm{\Lt-1}{\sigma} \Bm{\Lt-2}{\sigma} \cdots \Bm{0}{\sigma},
\qquad
\Bm{l}{\sigma}(s) = e^{-\dt K} e^{V_{l,\sigma}(s)},
\label{eq:Bdef}
\end{equation}
where \(V_{l,\sigma}(s) = \alpha\,\sigma\,\mathrm{diag}(s_{l,1},\ldots,s_{l,\Ns})\)
(here \(\sigma=\pm 1\) for \(\uparrow,\downarrow\)). The \(B\) matrices
are \(\Ns\times\Ns\) and are computed lazily, never stored, in the
implementation.

\paragraph{Sign-free regime.}
On bipartite lattices at \(\mu=0\) particle--hole symmetry maps
\(c_{i\sigma} \mapsto (-1)^{a_i} c^{\dagger}_{i\sigma}\) (with
\(a_i\) the sublattice label) and exchanges the two spin sectors,
giving \(\det M_\uparrow(s) = \det M_\downarrow(s)\) for every aux-field
configuration. The product \(\det M_\uparrow \det M_\downarrow \ge 0\),
so the Monte Carlo weight is non-negative and the sign problem is
absent. Away from half-filling the sign decays as
\(\langle\mathrm{sign}\rangle \sim e^{-\beta U f(\rho,T)}\) and rapidly
becomes the limiting factor in low-temperature simulations of the doped
Hubbard model.

\section{Equal-time Green's function}

The fundamental object measured by DQMC is the equal-time Green's
function on time slice \(l\),
\begin{equation}
G_\sigma(l)_{ij} = \big\langle c_{i\sigma}\, c^{\dagger}_{j\sigma} \big\rangle_{l}
                 = \big[\mathbb{1} + A_\sigma(l)\big]^{-1}_{ij},
\label{eq:Gdef}
\end{equation}
where the cyclic ordering of the propagator product is the
\emph{convention} that determines the form of the local-update
acceptance ratio. We adopt the \emph{pre-\(B_l\)} convention,
\begin{equation}
A_\sigma(l) =
  \Bm{l-1}{\sigma}\Bm{l-2}{\sigma}\cdots
  \Bm{0}{\sigma}\Bm{\Lt-1}{\sigma}\cdots\Bm{l+1}{\sigma}\Bm{l}{\sigma},
\label{eq:Aconv}
\end{equation}
which leads (Sec.~\ref{sec:update}) to the simple ratio
\(r = 1 + \Delta(1-G_{ii})\). All Wick-theorem-based observables are
expressed below in the same convention.

\section{Single-spin-flip update}\label{sec:update}

A flip \(s_{l,i} \mapsto -s_{l,i}\) modifies only \(\Bm{l}{\sigma}\):
\(\Bm{l}{\sigma}\to \Bm{l}{\sigma}\bigl(\mathbb{1} +
\Delta_\sigma E_{ii}\bigr)\) where \(E_{ii}\) is the rank-one matrix
with a single \(1\) at position \((i,i)\) and
\begin{equation}
\Delta_\uparrow   = e^{-2\alpha s_{l,i}} - 1,
\qquad
\Delta_\downarrow = e^{+2\alpha s_{l,i}} - 1.
\end{equation}
In the pre-\(B_l\) convention, this multiplies \(A\) on the
\emph{right}: \(A' = A(\mathbb{1} + \Delta E_{ii})\). A single Sherman--Morrison
identity then gives the determinant ratio
\begin{equation}
\boxed{\;
R_\sigma \equiv \frac{\det(\mathbb{1} + A'_\sigma)}{\det(\mathbb{1} + A_\sigma)}
       = 1 + \Delta_\sigma\bigl(1 - G_\sigma(i,i)\bigr).
\;}
\label{eq:Rratio}
\end{equation}
The Metropolis acceptance probability is
\(p = \min\!\bigl(1,\; \lvert R_\uparrow R_\downarrow\rvert\bigr)\).

If the flip is accepted, the equal-time Green's function is updated
exactly by another Sherman--Morrison step. Setting
\(u_j = \delta_{ji} - G_\sigma(j,i)\) and noting that
\(v = e_i\),
\begin{equation}
G'_\sigma = \biggl(\mathbb{1} - \frac{u\, e_i^{\!\top}}{R_\sigma}\biggr) G_\sigma,
\qquad
G'_\sigma(j,k) = G_\sigma(j,k)
   - \frac{\Delta_\sigma}{R_\sigma}
     \bigl(\delta_{ji} - G_\sigma(j,i)\bigr) G_\sigma(i,k).
\label{eq:SMupdate}
\end{equation}
This is an \(O(\Ns^2)\) outer product per accepted move (and
\(O(\Ns)\) per proposal for the ratio).

\section{Wrapping between time slices}

After all single-spin updates at slice \(l\) we advance to slice
\(l+1\). From the cyclic ordering of \(A(l)\) one gets directly
\begin{equation}
A_\sigma(l+1) = \Bm{l}{\sigma} \, A_\sigma(l) \, \Bm{l}{\sigma}^{-1},
\qquad
\boxed{\;
G_\sigma(l+1) = \Bm{l}{\sigma} \, G_\sigma(l) \, \Bm{l}{\sigma}^{-1}.
\;}
\label{eq:wrap}
\end{equation}
With \(\Bm{l}{\sigma} = e^{-\dt K} \cdot \mathrm{diag}(d)\) and
\(\Bm{l}{\sigma}^{-1} = \mathrm{diag}(d^{-1}) \cdot e^{+\dt K}\),
the order of operations matters because \(\mathrm{diag}(d^{-1})\) does
\emph{not} commute with \(e^{+\dt K}\):
\begin{equation}
G_\sigma(l+1)
  = e^{-\dt K} \,
    \bigl[\mathrm{diag}(d) \cdot G_\sigma(l) \cdot \mathrm{diag}(d^{-1})\bigr]
    \, e^{+\dt K} .
\label{eq:wrap-impl}
\end{equation}
\Cref{eq:wrap-impl} is what the code in \texttt{wrap\_green\_up\_} and
\texttt{wrap\_green\_dn\_} actually computes; getting this ordering
wrong is a subtle source of physical error that \emph{vanishes
identically} at \(U=0\), so the free-fermion regression test does not
catch it.

\section{Numerical stabilization}

The matrix product
\(A_\sigma(l) = \Bm{l-1}{\sigma}\cdots\Bm{l}{\sigma}\) is catastrophically
ill-conditioned at low temperatures: each \(e^{-\dt K}\) factor stretches
some directions by \(e^{+\dt \lambda_{\max}(K)}\) and contracts others by
\(e^{-\dt \lambda_{\min}(K)}\), so the singular values of \(A\) span
\(\sim e^{2\beta\lVert K\rVert}\) and the naive product loses all precision
within a few tens of slices.

We follow the standard remedy of Loh and
Gubernatis~\cite{LohGubernatis1992}: maintain an explicit
QR factorisation \((Q,R)\) of the partial product as we left-multiply by
each \(B\). Concretely:
\begin{enumerate}[leftmargin=*]
    \item Initialise \(Q = R = \mathbb{1}\) and an accumulator \(C = \mathbb{1}\).
    \item For \(\text{step} = 0, 1, \ldots, \Lt - 1\): set
          \(C \leftarrow B_{(l + \text{step})\!\!\mod \Lt}\, C\). Every
          \texttt{n\_stab} steps, refactorise
          \(C\,Q = Q' R'\) by Householder QR, then update
          \(Q \leftarrow Q'\), \(R \leftarrow R'\,R\), \(C \leftarrow \mathbb{1}\).
    \item After the last step, the partial product equals \(Q R\)
          exactly, with \(Q\) orthogonal to working precision.
\end{enumerate}
Then
\begin{equation}
G_\sigma = (\mathbb{1} + Q R)^{-1}
        = \bigl(Q\,(Q^{\!\top} + R)\bigr)^{-1}
        = (Q^{\!\top} + R)^{-1}\, Q^{\!\top},
\label{eq:Gqr}
\end{equation}
which is solved by a plain LU on \((Q^{\!\top} + R)\). Even though
\(R\) itself can have a huge dynamical range, the linear system
\eqref{eq:Gqr} is conditioned by \(G\) (whose entries are
\(O(1)\)), and the solve is stable. In our test runs the
\texttt{max\_recompute\_drift} reported by the engine -- the
\(\ell^\infty\) discrepancy between the wrapped \(G\) and the freshly
recomputed \(G\) -- stays at machine precision (\(\le 10^{-12}\)) for
all parameter regimes we have tried up to \(\beta U \sim 32\),
\(\Ns = 36\).

\section{Observables}

For a fixed Hubbard--Stratonovich configuration the system is
non-interacting and Wick's theorem applies exactly. With our convention
\(G_\sigma(i,j)=\langle c_{i\sigma} c^{\dagger}_{j\sigma}\rangle\) we have
\(\langle n_{i\sigma}\rangle = 1 - G_\sigma(i,i)\) and
\begin{equation}
\langle S^z_i S^z_j \rangle = \tfrac{1}{4}\Bigl[
  (1 - G_\uparrow(i,i))(1 - G_\uparrow(j,j))
+ (1 - G_\downarrow(i,i))(1 - G_\downarrow(j,j))
\end{equation}
\vspace{-2em}
\begin{equation}
\hphantom{\langle S^z_i S^z_j \rangle = \tfrac{1}{4}\Bigl[}
- 2(1 - G_\uparrow(i,i))(1 - G_\downarrow(j,j))
+ (\delta_{ij} - G_\uparrow(j,i))\,G_\uparrow(i,j)
+ (\uparrow\!\leftrightarrow\!\downarrow)
\Bigr].
\end{equation}
The implementation evaluates
\(\langle n\rangle\), the double occupancy
\(\frac{1}{\Ns}\sum_i (1-G_\uparrow(i,i))(1-G_\downarrow(i,i))\),
\(\langle m_z^2\rangle = (1/\Ns^2)\sum_{ij}\langle S^z_iS^z_j\rangle\),
the staggered moment
\(\langle m_{\text{stag}}^2\rangle = (1/\Ns^2)\sum_{ij} (-1)^{a_i+a_j}
   \langle S^z_iS^z_j\rangle\),
and the antiferromagnetic structure factor
\(S(\pi,\pi) = (1/\Ns)\sum_{ij}(-1)^{a_i+a_j}\langle S^z_iS^z_j\rangle\).
Errors use the same Flyvbjerg--Petersen logarithmic binning analysis as
the SSE code (\texttt{qmc::Observable}).

\section{Verification}

The free-fermion limit ($U=0$) is the most powerful regression test:
the auxiliary field decouples, every configuration produces the same
analytical Green's function, and the simulation must reproduce it to
machine precision. Concretely,
\begin{equation}
\bigl\lvert G_\sigma^{\text{DQMC}} - (\mathbb{1} + e^{-\beta K})^{-1} \bigr\rvert_{\infty}
   < 10^{-13}
\end{equation}
both at initialisation and after a full sweep on the L=8 chain at
$\beta = 4$ (\texttt{tests/test\_dqmc.cpp}).

A second exact identity ---
\(\langle n\rangle = 1\) per configuration on a half-filled bipartite
lattice --- is reproduced to $\sim 10^{-9}$ on the L=4 chain at
$U=4,\beta=2$. The reported average sign is exactly $+1$ in this
sign-free regime.

\Cref{tab:dqmc-bench} summarises a handful of moderate runs that
illustrate the expected physics:
\begin{table}[h!]
\centering
\begin{tabular}{lccccc}
\toprule
System & $\beta$ & $\langle n\rangle$ & $\langle n_\uparrow n_\downarrow\rangle$
       & $S(\pi,\pi)$ & $\langle\mathrm{sign}\rangle$\\
\midrule
chain $L=8$, $U=2$  & 4 & $1.000$ & $0.169$ & $0.36$  & $1.0$\\
square $4{\times}4$, $U=4$  & 4 & $1.000$ & $0.124$ & $0.66$  & $1.0$\\
square $6{\times}6$, $U=4$  & 8 & $1.000$ & $0.120$ & $1.58$  & $1.0$\\
\bottomrule
\end{tabular}
\caption{Equal-time observables on three half-filled benchmarks. The
growth of $S(\pi,\pi)$ with system size and $\beta$ reflects the
expected antiferromagnetic ordering tendency of the half-filled square
Hubbard model.}
\label{tab:dqmc-bench}
\end{table}

\section*{Code map}

\begin{table}[h!]
\centering
\small
\begin{tabular}{ll}
\toprule
File & Purpose\\
\midrule
\texttt{include/qmc/linalg.hpp}            & Dense matrix, GEMM, LU, QR, exp$(-\dt K)$, stable $G$\\
\texttt{include/qmc/hubbard.hpp}           & Hubbard parameters, kinetic matrix\\
\texttt{include/qmc/dqmc\_engine.hpp}      & B-builder, sweep, Sherman--Morrison, wrap, recompute\\
\texttt{include/qmc/dqmc\_measurements.hpp}& Equal-time observables (Wick contractions)\\
\texttt{apps/qmc\_dqmc.cpp}                & CLI driver, OpenMP parallel replicas\\
\texttt{tests/test\_dqmc.cpp}              & Free-fermion + half-filling regression tests\\
\texttt{tests/test\_linalg.cpp}            & Linear algebra unit tests\\
\bottomrule
\end{tabular}
\end{table}

\begin{thebibliography}{9}
\bibitem{BSS1981}
R.~Blankenbecler, D.~J.~Scalapino, R.~L.~Sugar,
\emph{Monte Carlo calculations of coupled boson--fermion systems. I},
Phys. Rev. D \textbf{24}, 2278 (1981).

\bibitem{Hirsch1983}
J.~E.~Hirsch,
\emph{Discrete Hubbard--Stratonovich transformation for fermion lattice models},
Phys. Rev. B \textbf{28}, 4059 (1983).

\bibitem{LohGubernatis1992}
E.~Y.~Loh~Jr.\ and J.~E.~Gubernatis,
\emph{Stable numerical simulations of models of interacting electrons in
condensed-matter physics}, in
\emph{Electronic Phase Transitions}, ed.\ W.~Hanke and Yu.~V.~Kopaev,
Modern Problems in Condensed Matter Sciences \textbf{32}, 177 (1992).

\bibitem{Assaad2008}
F.~F.~Assaad and H.~G.~Evertz,
\emph{World-line and Determinantal Quantum Monte Carlo Methods for
Spins, Phonons and Electrons},
Lect.\ Notes Phys.\ \textbf{739}, 277 (2008).
\end{thebibliography}

\end{document}
